\section{Discussion}
\label{sec:discussion}

\subsection{Limitations}

\paragraph{SCM structure learning as channel optimization, not
prerequisite.} The causal attribution channel of
Section~\ref{sec:causal} requires the actor to know or learn the
causal DAG over agent actions and capability changes. Standard
observational structure learning \citep{spirtes2000causation} is
inadequate for the GFM setting for three distinct reasons. First,
the data-generating process is adversarial: scorpions can
strategically align their actions with observed noise to mask their
causal role, so classical conditional-independence tests can be
gamed by an adversary who knows which tests will be run. Standard
structure learning does not have a threat model. Second, the agent
cardinality is non-fixed: new agents can appear, existing agents can
delegate, and a group can coordinate and behave as a single
effective actor, so the variable set itself evolves. Latent-variable
discovery methods (FCI, RFCI) assume a fixed variable cardinality
and do not naturally handle this. Third, the actor cannot typically
intervene on other agents: interventional causal discovery relies on
actual interventions, which a GFM actor can simulate inside its SCM
but cannot enforce in the world. This rules out the entire class of
interventional discovery methods and leaves only observational
methods, which are strictly weaker under the adversarial conditions
above.

Because causal attribution is one channel among several
(Section~\ref{sec:multi_channel}), and capability gating via the
poset of \citet{lasser2026poset} and \citet{lasser2026horizon} is
the strongest and most difficult to fake channel (relative to a
trusted, sufficiently complete capability graph), structure learning
becomes an \emph{optimization target} for one channel's contribution
rather than a \emph{prerequisite} for attribution at all. When the
SCM is poorly learned, the causal channel degrades, the other
channels are not directly degraded by the SCM misspecification, and
the multi-channel aggregate reweights toward the
structurally-grounded channels. The system does not collapse; it
falls back on the trust and capability machinery of
\citet{lasser2026gfm}, \citet{lasser2026poset}, and
\citet{lasser2026horizon}.

The structural assumptions underlying
Proposition~\ref{prop:causal_dominates}(c) are partially testable
through a residual-covariance diagnostic: if
$\mathrm{Cov}(\CausalAttr(a_j), \CausalAttr(a_k) \mid G_t)$ exceeds
the level the current SCM predicts, the SCM is misspecified ---
possible causes include a missing edge, latent common causes,
exogenous-orthogonality failure, non-additivity, or nonstationarity
--- and the causal channel's reliability should degrade accordingly. Integrating such a
diagnostic into a calibrated SCM-confidence scalar
$c_{\SCM}(t) \in [0,1]$ that weights the causal channel's
contribution to the multi-channel aggregate --- so declining
sufficiency down-weights causal attribution without burning specific
agents' trust --- is an architectural direction we identify but do
not formalize here. The full adversarial-robust structure-learning
problem, including the calibration procedure and its interaction
with the behavioral and capability channels, is deferred to a
companion paper.

\paragraph{Risk-trust residual dependence on the actor's model.} The
pathway-conditional structural verification residual
(Definition~\ref{def:risk_residual}) measures divergence between the
claim and the actor's own assessment of the claimed pathway. If the
actor's world model is itself poorly calibrated, the residual reflects
model error rather than claim quality. An agent whose risk assessments
are accurate but disagree with a miscalibrated actor will accumulate
high $\|r^{\mathrm{risk}}\|^2$ and concentrate around low
$\Trisk_j^*$ --- a false negative. This is mitigated by the self-trust
mechanism of \citet{lasser2026gfm}: an actor with a poorly calibrated
model has low $T_s$ and high learning rate, which accelerates model
correction. But during the correction period, risk-trust assessments
may be unreliable.

\paragraph{Log-odds aggregation and reporter correlation.} The
trust-weighted log-linear opinion pool (Definition~\ref{def:logodds})
recovers Bayesian updating in the no-tempering limit under conditional
independence of agents' observations given the true risk level. If
reporters' assessments are correlated --- they share information
sources, or one reporter's assessment influences others --- the
aggregation overweights the shared signal. In practice, partial
independence is the norm: agents with different observation channels
and different world models produce partially independent assessments,
and the log-linear pool is an approximation whose error grows with
the correlation. Characterizing the approximation error under partial
independence, and deriving a correction term from a reporter-network
model, is deferred to future work.

